\documentclass[aps,preprint]{revtex4}%
\usepackage{amsfonts}
\usepackage{amsmath}
\usepackage{amssymb}
\usepackage{graphicx}%
\setcounter{MaxMatrixCols}{30}
%TCIDATA{OutputFilter=latex2.dll}
%TCIDATA{Version=4.10.0.2363}
%TCIDATA{CSTFile=revtex4.cst}
%TCIDATA{Created=Monday, September 20, 2004 09:58:21}
%TCIDATA{LastRevised=Saturday, November 20, 2004 09:34:41}
%TCIDATA{<META NAME="GraphicsSave" CONTENT="32">}
%TCIDATA{<META NAME="DocumentShell" CONTENT="Articles\SW\REVTeX 4">}
\newtheorem{theorem}{Theorem}
\newtheorem{acknowledgement}[theorem]{Acknowledgement}
\newtheorem{algorithm}[theorem]{Algorithm}
\newtheorem{axiom}[theorem]{Axiom}
\newtheorem{claim}[theorem]{Claim}
\newtheorem{conclusion}[theorem]{Conclusion}
\newtheorem{condition}[theorem]{Condition}
\newtheorem{conjecture}[theorem]{Conjecture}
\newtheorem{corollary}[theorem]{Corollary}
\newtheorem{criterion}[theorem]{Criterion}
\newtheorem{definition}[theorem]{Definition}
\newtheorem{example}[theorem]{Example}
\newtheorem{exercise}[theorem]{Exercise}
\newtheorem{lemma}[theorem]{Lemma}
\newtheorem{notation}[theorem]{Notation}
\newtheorem{problem}[theorem]{Problem}
\newtheorem{proposition}[theorem]{Proposition}
\newtheorem{remark}[theorem]{Remark}
\newtheorem{solution}[theorem]{Solution}
\newtheorem{summary}[theorem]{Summary}
\newenvironment{proof}[1][Proof]{\noindent\textbf{#1.} }{\ \rule{0.5em}{0.5em}}
\begin{document}
\preprint{HEP/123-qed}
\title[Short title for running header]{Long title that appears on the title page}
\author{First Author}
\affiliation{My Institution}
\author{Second Author}
\affiliation{My Institution}
\author{Third Author}
\affiliation{Other Institution}
\keywords{one two three}
\pacs{PACS number}

\begin{abstract}
Shell document for REV\TeX{} 4.

\end{abstract}
\volumeyear{year}
\volumenumber{number}
\issuenumber{number}
\eid{identifier}
\date[Date text]{date}
\received[Received text]{date}

\revised[Revised text]{date}

\accepted[Accepted text]{date}

\published[Published text]{date}

\startpage{101}
\endpage{102}
\maketitle
\tableofcontents


\section{Introduction}

The total energy can be expressed as
\begin{equation}
\mathcal{E}\left(  \mathbb{D}^{2}\left(  \bm{\upsilon}\right)
,\bm{\upsilon}\right)  =\mathcal{E}\left(  \mathbb{D}^{2}\left(
\bm{\varphi}\right)  ,\bm{\varphi}\right)  =\frac{1}{2}\sum_{1\leq p\leq
2s}N_{p}\left\{  I\left(  \upsilon_{p}\right)  +\left\langle \left.
\upsilon_{p}\right\vert h_{1}\upsilon_{p}\right\rangle \right\}  =\frac{1}{2}%
%TCIMACRO{\tsum \limits_{1\leq p\leq2s}}%
%BeginExpansion
{\textstyle\sum\limits_{1\leq p\leq2s}}
%EndExpansion
N_{p}\kappa\left(  \upsilon_{p}\right)
\end{equation}
where
\begin{equation}
\kappa\left(  \upsilon_{p}\right)  =I\left(  \upsilon_{p}\right)
+\left\langle \left.  \upsilon_{p}\right\vert h_{1}\upsilon_{p}\right\rangle
\end{equation}%
\begin{equation}
h_{1}=\sum_{1\leq j,k\leq2s}h_{1jk}\left(  \upsilon\right)  \left\vert
\upsilon_{j}\right\rangle \left\langle \upsilon_{k}\right\vert
\end{equation}
and
\begin{equation}
h_{1jk}\left(  \upsilon\right)  =\left\langle \upsilon_{j}\left\vert h\right.
\upsilon_{k}\right\rangle =-\int\upsilon_{j}^{\ast}\left(  \bm{r,\xi}\right)
\left\{  \sum_{_{1\leq l\leq N_{nuc}}}\frac{Z_{l}}{\left\Vert \bm{R}_{l}%
-\bm{r}\right\Vert }+\nabla_{\bm{r}}^{2}\right\}  \upsilon_{k}\left(
\bm{r,\xi}\right)  d\bm{r}d\xi;\quad1\leq j,k\leq2s,\label{h1}%
\end{equation}
Using the Fock operator
\begin{equation}
\mathcal{E}\left(  \mathbb{D}^{2},\bm{\varphi}\right)  =\frac{1}{2}\left\{
\overset{}{\operatorname{Tr}\left\{  F\left(  \mathbb{D}^{2}%
,\bm{\varphi}\right)  \right\}  +\mathcal{E}_{1}\left(  \mathbb{D}%
^{2},\bm{\varphi}\right)  }\right\}  .
\end{equation}%
\begin{align}
\mathcal{E}\left(  \mathbf{n,\psi}\right)   &  =\frac{\mathcal{H}\left(
\mathbf{n,\psi}\right)  }{S_{N}\left(  \mathbf{n}\right)  }\nonumber\\
&  =\frac{1}{S_{N}\left(  \mathbf{n}\right)  }\left\{  \sum_{1\leq j\leq
s}n_{j}S_{N-1}\left(  \jmath\right)  \left\{  \left\langle \psi_{j}\left\vert
h_{1}\psi_{j}\right.  \right\rangle +\left\langle \psi_{j+s}\left\vert
h_{1}\psi_{j+s}\right.  \right\rangle +\left\langle \psi_{j}\psi_{j+s}%
||\psi_{j}\psi_{j+s}\right\rangle \right\}  \right.  \nonumber\\
&  +\sum_{1\leq j,k\leq s}n_{j}^{\frac{1}{2}}n_{k}^{\frac{1}{2}}S_{N-1}\left(
\jmath k\right)  \left\langle \psi_{j}\psi_{j+s}||\psi_{k}\psi_{k+s}%
\right\rangle \left.  +\sum_{1\leq j<k\leq s}n_{i}n_{j}S_{N-2}\left(  \jmath
k\right)  \left\langle \psi_{j}\psi_{k}|||\psi_{j}\psi_{k}\right\rangle
\right\}  .\label{energy}%
\end{align}%
\begin{equation}
\mathcal{E}\left(  \mathbf{\eta},\mathbf{\lambda}\right)  =\frac
{\mathcal{H}\left(  \mathbf{\eta},\mathbf{\lambda}\right)  }{S_{N}\left(
\mathbf{\eta}\right)  }=\frac{\left\langle \left.  g\left(  \mathbf{\eta
},\mathbf{\lambda}\right)  ^{N}\right\vert Hg\left(  \mathbf{\eta
},\mathbf{\lambda}\right)  ^{N}\right\rangle }{S_{N}\left(  \mathbf{\eta
}\right)  },\label{gp_funct}%
\end{equation}%
\begin{equation}
S_{N}\left(  \mathbf{\eta}\right)  =\left\langle \left.  k\left(
\mathbf{\eta}\right)  g_{0}^{N}\right\vert k\left(  \mathbf{\eta}\right)
g_{0}^{N}\right\rangle =\sum_{1\leq j_{1}<\cdots<j_{N}\leq s}e^{\eta_{j_{1}%
}+\cdots+\eta_{j_{N}}}n_{j_{1}}\cdots n_{j_{N}}.
\end{equation}%
\begin{align}
\mathcal{L}\left(  \mathbf{n,\psi}\right)   &  =\mathcal{E}\left(
\mathbf{n,\psi}\right)  -\sum_{1\leq j,k\leq2s}\varepsilon_{kj}\left(
\left\langle \left.  \psi_{j}\right\vert \psi_{k}\right\rangle -\delta
_{jk}\right)  \label{larE}\\
\frac{\partial\mathcal{L}\left(  \mathbf{n,\psi}\right)  }{\partial n_{j}} &
=\frac{\partial\mathcal{E}\left(  \mathbf{n,\psi}\right)  }{\partial n_{j}}.
\end{align}%
\begin{align}
\mathcal{L}\left(  \mathbf{\psi,s}\right)   &  =\mathcal{E}\left(
\mathbf{n}_{0}\mathbf{,\psi}\right)  -\sum_{1\leq j,k\leq2s}\varepsilon
_{kj}\left(  \left\langle \left.  \psi_{j}\right\vert \psi_{k}\right\rangle
-s_{j}\delta_{jk}\right)  \\
\frac{\partial\mathcal{L}\left(  \mathbf{\psi,s}\right)  }{\partial s_{j}} &
=\varepsilon_{jj}.
\end{align}%
\begin{equation}
\mathcal{L}\left(  \mathbf{n\left(  \mathbf{\eta}\right)  ,\psi}\right)
=\mathcal{L}\left(  \mathbf{\psi,s}\left(  \mathbf{\eta}\right)  \right)
\end{equation}
The occupation numbers of the AGP\ state are
\begin{equation}
N_{j}=\mathcal{N}_{j}\left(  \mathbf{n}\right)  =n_{j}\frac{S_{N-1}\left(
\jmath\right)  }{S_{N}}=n_{j}\frac{\partial\ln S_{N}}{\partial n_{j}}%
;\quad1\leq j\leq s.
\end{equation}
The phase variation \emph{unitary} subgroup $U_{phase}$ is given by%
\begin{equation}
U_{phase}=\left\{  \left.  p\left(  \bm{\theta}\right)  =\exp\left\{
i\sum_{1\leq j\leq s}\theta_{j}\Omega_{j}\right\}  \right\vert 0\leq\theta
_{j}<2\pi\right\}  ,\label{phase}%
\end{equation}
where
\begin{equation}
\Omega_{j}=\left(  a_{j}^{\dagger}a_{j}+a_{j+s}^{\dagger}a_{j+s}\right)
;\quad1\leq j\leq s
\end{equation}
and occupation number variation nonunitary subgroup $G_{num}$ by
\begin{equation}
G_{num}=\left\{  \left.  k\left(  \bm{\eta}\right)  =\exp\left\{  \sum_{1\leq
j\leq s}\eta_{j}\Omega_{j}\right\}  \right\vert \eta_{j}\in\mathbb{R}\right\}
.\label{occup}%
\end{equation}
The manifold of $2N$-electron AGP\ states $\mathcal{M}$ can then be expressed
as
\begin{align}
\mathcal{M} &  \mathcal{=}\left\{  \left.  \overset{}{k\left(
\bm{\eta}\right)  p\left(  \bm{\theta}\right)  u\left(  \bm{\alpha}\right)
}\left\vert g_{0}^{N}\right\rangle \right\vert \quad0\leq\theta_{j}%
<2\pi;\,\eta_{j}\in\mathbb{R};\right.  \nonumber\\
&  \qquad\qquad\qquad\qquad\qquad\left.  \bm{\alpha}=\bm{\alpha }^{\dagger
},\operatorname{Tr}\bm{\alpha}=0,1\leq j\neq k\pm s\leq2s\in\mathbb{C}%
\right\}
\end{align}
Explicitly the action on the reference geminal is given by
\begin{equation}
\left\vert g\left(  \bm{\eta},\bm{\alpha},\bm{\theta}\right)  \right\rangle
=k\left(  \bm{\eta}\right)  p\left(  \bm{\theta}\right)  u\left(
\bm{\alpha}\right)  \left\vert g_{0}\right\rangle =\sum_{1\leq l\leq s}%
e^{\eta_{l}}e^{i\theta_{l}}c_{0l}\left\vert \varphi_{l}\left(
\bm{\alpha}\right)  \varphi_{l+s}\left(  \bm{\alpha}\right)  \right\rangle ,
\end{equation}
where
\begin{equation}
\left\vert \varphi_{l}\left(  \bm{\alpha}\right)  \right\rangle =\exp\left\{
i\sum_{1\leq j\neq k\pm s\leq2s}\alpha_{jk}a_{j}^{\dagger}a_{k}\right\}
\left\vert \varphi_{0l}\right\rangle .
\end{equation}
One can define a normalized geminal by
\begin{equation}
\end{equation}
The normalization condition
\begin{equation}
\left\langle \left.  g\left(  \bm{\eta},\bm{\alpha},\bm{\theta}\right)
\right\vert g\left(  \bm{\eta},\bm{\alpha},\bm{\theta}\right)  \right\rangle
=1
\end{equation}
gives
\begin{equation}
\sum_{1\leq l\leq s}e^{2\eta_{l}}n_{l}%
\end{equation}
when We can define%
\begin{equation}
\left\vert \varphi_{l}\left(  \bm{\eta,\theta,\alpha}\right)  \right\rangle
=\exp\frac{\eta_{l}}{2}\exp i\frac{\theta_{l}}{2}\exp\left\{  i\sum_{1\leq
j\neq k\pm s\leq2s}\alpha_{jk}a_{j}^{\dagger}a_{k}\right\}  \left\vert
\varphi_{0l}\right\rangle ,
\end{equation}
so that
\begin{equation}
\left\langle \left.  \varphi_{l}\left(  \bm{\eta,\theta,\alpha}\right)
\right\vert \varphi_{l}\left(  \bm{\eta,\theta,\alpha}\right)  \right\rangle
=e^{\eta_{l}}\left\langle \left.  \varphi_{0l}\right\vert \varphi
_{0l}\right\rangle .
\end{equation}



\end{document}